\subsection{LIE ALGEBRAS}

DEFINITION 1. An algebra g over K is called a Lie algebra over K if its multiplication (denoted by \((x,y)\mapsto [x,y])\) satisfies the identities:

\[
[ x, x ] = 0\tag{1}
\]

\[
[ x, [ y, z ] ] + [ y, [ z, x ] ] + [ z, [ x, y ] ] = 0\tag{2}
\]

for all x, y, z in g.

The product \([x, y]\) is called the bracket of \(x\) and \(y\) . Identity (2) is called the Jacobi identity.

The bracket \([x, y]\) is an alternating bilinear function of \(x\) and \(y\) . We have the identity:

\[
[ x, y ] = - [ y, x ]\tag{3}
\]

so that the Jacobi identity can be written:

\[
[ x, [ y, z ] ] = [ [ x, y ], z ] + [ y, [ x, z ] ].\tag{4}
\]

Every subalgebra and every quotient algebra of a Lie algebra is a Lie algebra. Every product of Lie algebras is a Lie algebra. If g is a Lie algebra, the opposite algebra \(g^{0}\) is a Lie algebra and the mapping \(x \mapsto -x\) an isomorphism of g onto \(g^{0}\) , by virtue of identity (3).

Example 1. Let L be an associative algebra over K. The bracket \([x, y] = xy - yx\) is a bilinear function of \(x\) and \(y\) . It is easily verified that the law of composition \((x, y) \mapsto [x, y]\) on the K-module L makes L into a Lie algebra over K.

Example 2. In Example 1 choose L to be the associative algebra of endomorphisms of a K-module E. We obtain the Lie algebra of endomorphisms of E, denoted by \(\mathfrak{gl}(\mathbf{E})\) . (If \(\mathbf{E} = \mathbf{K}^n\) , the Lie algebra \(\mathfrak{gl}(\mathbf{E})\) is denoted by \(\mathfrak{gl}(n, \mathbf{K})\) .)

Every Lie subalgebra of \(\mathfrak{gl}(\mathbf{E})\) is a Lie algebra over K. In particular:

\begin{enumerate}
\def\labelenumi{(\arabic{enumi})}
\item
  If E is given a (not necessarily associative) algebra structure, the derivations of E form a Lie algebra over K.
\item
  If E admits a finite basis, the endomorphisms of E of zero trace form a Lie algebra over K denoted by \(\mathfrak{sl}(\mathbf{E})\) (or \(\mathfrak{sl}(n,\mathbf{K})\) if \(E = K^{n}\) ).
\item
  The set \(\mathbf{M}_n(\mathbf{K})\) of square matrices of order \(n\) can be considered as a Lie algebra over K canonically isomorphic to \(\mathfrak{gl}(n, \mathbf{K})\) . Let \((E_{ij})\) be the canonical basis of \(\mathbf{M}_n(\mathbf{K})\) (Algebra, Chapter II, § 10, no. 3). It follows easily that:
\end{enumerate}

\[
\left\{ \begin{array}{l l} [ E _ {i j}, E _ {k l} ] = 0 & \text {if} j \neq k \quad \text {and} \quad i \neq l \\ {} [ E _ {i j}, E _ {j l} ] = E _ {i l} & \text {if} i \neq l \\ {} [ E _ {i j}, E _ {k i} ] = - E _ {k j} & \text {if} j \neq k \\ {} [ E _ {i j}, E _ {j i} ] = E _ {i i} - E _ {j j} \end{array} \right.\tag{5}
\]

The Lie subalgebra of \(\mathbf{M}_{n}(\mathbf{K})\) consisting of the triangular matrices (resp. triangular matrices of zero trace, resp. triangular matrices of zero diagonal) is denoted by \(t(n,\mathbf{K})\) (resp. \(\mathfrak{st}(n,\mathbf{K})\) , resp. \(\mathfrak{n}(n,\mathbf{K})\) ) (Algebra, Chapter II, § 10, no. 7).

*Example 3. Let V be an infinitely differentiable real manifold. The differential operators with infinitely differentiable real coefficients constitute an associative algebra over R and hence, by Example 1, a Lie algebra \(\Delta\) over R. The bracket of two infinitely differentiable vector fields on V is an infinitely differentiable vector field and hence the infinitely differentiable vector fields on V constitute a Lie subalgebra \(\mathfrak{f}\) of \(\Delta\) . If V is a real Lie group, the left invariant vector fields constitute a Lie subalgebra \(\mathfrak{g}\) of \(\mathfrak{f}\) called the Lie algebra of V. The vector space \(\mathfrak{g}\) is identified with the tangent space to V at \(e\) (the identity element of V). Let V' be another real Lie group, \(e'\) its identity element and \(\mathfrak{g}'\) its Lie algebra. Every analytic homomorphism of V into V' defines a linear mapping of the tangent space to V at \(e\) into the tangent space to V' at \(e'\) ; this mapping is a homomorphism of the Lie algebra \(\mathfrak{g}\) into the Lie algebra \(\mathfrak{g}'\) . If V is the linear group of a finite-dimensional real vector space E there exists a canonical isomorphism of gl(E) onto the Lie algebra \(\mathfrak{g}\) of V, under which \(\mathfrak{g}\) is identified with gl(E) \(_*\) .

DEFINITION 2. Let g be a Lie algebra and x an element of g. The linear mapping \(y \mapsto [x, y]\) of g into g is called the adjoint linear mapping of x and is denoted by \(ad_{g}x\) or ad x.

PROPOSITION 1. Let \(\mathfrak{g}\) be a Lie algebra. For all \(x \in \mathfrak{g}\) , ad \(x\) is a derivation. The mapping \(x \mapsto \operatorname{ad} x\) is a homomorphism of the Lie algebra \(\mathfrak{g}\) into the Lie algebra \(\mathfrak{d}\) of derivations of \(\mathfrak{g}\) . If \(\mathrm{D} \in \mathfrak{d}\) and \(x \in \mathfrak{g}\) , {[}D, ad \(x\) {]} = \(\operatorname{ad}(\mathrm{D}x)\) .

Identity (4) can be written:

\[
(\operatorname{ad} x). [ y, z ] = [ (\operatorname{ad} x). y, z ] + [ y, (\operatorname{ad} x). z ]
\]

or:

\[
(\operatorname{ad} [ x, y ]). z = (\operatorname{ad} x). ((\operatorname{ad} y). z) - (\operatorname{ad} y). ((\operatorname{ad} x). z)
\]

whence the first two assertions. On the other hand, if \(\mathbf{D} \in \mathfrak{d}\) , \(x \in \mathfrak{g}\) , \(y \in \mathfrak{g}\) , then \([\mathbf{D}, \operatorname{ad} x] \cdot y = \mathrm{D}([x, y]) - [x, \mathrm{D}y] = [\mathrm{D}x, y] = (\operatorname{ad} \mathrm{D}x) \cdot y\) , whence the last assertion.

The mapping ad \(x\) is also called the inner derivation defined by \(x\) .

